\documentclass[10pt,a4paper]{article}
\usepackage[T1]{fontenc}
\usepackage[margin=1in]{geometry}
\usepackage{times}
\usepackage{xcolor}
\usepackage{enumitem}
\usepackage{booktabs}
\usepackage{parskip}

\definecolor{quotegray}{gray}{0.30}
\newenvironment{rev}%
  {\begin{quote}\itshape\color{quotegray}}%
  {\end{quote}}
\newcommand{\pt}[1]{\textbf{#1}}

\title{\vspace{-1.5em}Response to Reviewers\vspace{-0.5em}}
\author{}
\date{}

\begin{document}
\maketitle
\vspace{-3em}

\noindent\rule{\textwidth}{0.8pt}

\textbf{Paper:} How Much Does the Activation Function Matter?
A 2{,}000-Run Variance Study of Gated Recurrent Networks for
Sentiment Classification

\textbf{Conference:} iCONEECT 2026 \quad\textbf{Paper ID:} [PID]

\noindent\rule{\textwidth}{0.8pt}

We thank both reviewers. Reviewer~2 and Reviewer~3 independently
identified the same weakness, that the sweep was single-seed.
Following that point changed the paper substantially.

We did not restrict the replication to the top ten configurations
as asked. We re-ran \pt{every one of the 400 configurations under
five seeds, for 2{,}000 trained models}, because a partial
replication could not have told us whether the problem was local to
the leaders. It was not. The replication showed that the
single-seed ranking of activation pairs does not reproduce in any
of the four architectures, which we now report as the paper's
principal finding. Several claims in the submitted version did not
survive and have been removed or reversed; these are listed
explicitly at the end of this letter.

The title, abstract, research questions, results and conclusion
have been revised accordingly. The experimental design,
architectures, hyperparameters and reference list are unchanged.

\section*{Reviewer 2}

\begin{rev}
Methodology sound. Primary reason for revision: address single-seed
statistical variance in the results section.
\end{rev}

\noindent\pt{Addressed. See Section~III-E (Replication Protocol),
Section~IV-A, Table~IV, and Table~VI.}

Every configuration was retrained under five seeds
$\{10,20,30,40,50\}$ with the data split held fixed at
\texttt{random\_state=42}, so the reported variance is the variance
of retraining the same model on the same data. That is the quantity
a single-run comparison implicitly assumes to be negligible.

A variance decomposition over the 2{,}000 runs (Table~IV) gives:

\begin{center}
\begin{tabular}{lr}
\toprule
Source & \% of total variance \\
\midrule
Random seed & \pt{58.67} \\
Architecture & 26.69 \\
\quad of which cell type (GRU vs.\ LSTM) & 26.32 \\
\quad of which layer-2 bidirectionality & 0.19 \\
Act$_1\times$Act$_2$ pair & \pt{3.74} \\
\bottomrule
\end{tabular}
\end{center}

The seed accounts for more variance than every controlled factor
combined, and roughly sixteen times more than the activation pair.
Two conclusions changed as a result:

\begin{itemize}[leftmargin=1.4em,itemsep=2pt]
  \item \pt{The top-ten ranking is withdrawn.} The ten best of 400
  configurations span 0.12 accuracy points against confidence
  half-intervals up to $\pm$0.50. Every pair of intervals overlaps.
  Table~V now reports them as statistically indistinguishable and
  the text states that the ordering must not be read as a ranking.

  \item \pt{The architecture conclusion is strengthened.} GRU
  exceeds LSTM by 0.79 points over 1{,}000 runs per family
  ($p\approx2\times10^{-125}$, Cohen's $d=1.19$), up from the 0.73
  points reported from single runs. We additionally report that GRU
  is the more \emph{reliable} family, with roughly 40\% lower
  seed-to-seed standard deviation ($p\approx2\times10^{-7}$), and
  that all eight high-variance cells in the grid are LSTM-based.
\end{itemize}

We also added a power analysis (Section~V-F) giving the seed budget
needed to resolve each effect: about 4 seeds per cell for the
architecture difference, about 157 to separate the best cell from
the tenth. This makes the limitation quantitative rather than
rhetorical.

\section*{Reviewer 3}

\subsection*{Point 1: Figure quality}

\begin{rev}
Regenerate Fig.~1 and Fig.~2 as vector (PDF/EPS) or $\geq$300 DPI
raster. All axis titles, ticks, bar labels fully legible, no
clipping.
\end{rev}

\noindent\pt{Addressed. Both figures regenerated as vector PDF.}

Both are now produced by a script that embeds TrueType subsets
(\texttt{pdf.fonttype~=~42}), reads the trimmed bounding box back
out of the written PDF, and verifies the on-page size of the
smallest text. Fig.~1 renders at 8.6~pt and Fig.~2 at 8.3~pt, above
the 8~pt floor, with no clipped labels and legends placed clear of
the data.

Fig.~2 (formerly Fig.~1) is the convergence plot; the bar chart that
was Fig.~2 has been \pt{replaced rather than redrawn}. The submitted
version plotted the single best activation pair per architecture;
those peaks are maxima of a noisy single-seed sweep and do not
reproduce, so redrawing them would have made a misleading figure
merely legible. The new Fig.~1 plots all 400 replicated cell means
grouped by architecture with 95\% confidence intervals, which shows
directly what the data supports: the architectures separate, the
activation pairs within them do not.

Fig.~2 now shows the mean validation-loss trajectory over five
seeds with $\pm1$ standard deviation bands, rather than a single
run.

\subsection*{Point 2: Error bars and confidence intervals}

\begin{rev}
Main sweep is single-seed. Add standard deviation error bars or
confidence intervals for the top-10 configurations across multiple
seeds to show differences are statistically meaningful.
\end{rev}

\noindent\pt{Addressed, and extended to all 400 configurations.
See Table~V, Table~VI and Fig.~1.}

Table~V reports the ten best configurations with mean, standard
deviation and 95\% confidence interval, as requested. The direct
answer to the reviewer's question is that the differences are
\pt{not} statistically meaningful: the ten means span 0.12 points
and all intervals overlap.

Because we replicated the whole grid rather than only the leaders,
we can also report how far the single-seed ranking is from the
replicated one (Table~VI):

\begin{center}
\begin{tabular}{lccc}
\toprule
Experiment & Spearman $\rho$ & Mean rank $\Delta$ & Top-10 overlap \\
\midrule
BiLSTM+LSTM   & $-0.174$ & 36.9 & 0/10 \\
BiGRU+GRU     & $-0.059$ & 34.5 & 1/10 \\
BiLSTM+BiLSTM & $-0.041$ & 34.4 & 1/10 \\
BiGRU+BiGRU   & $+0.055$ & 32.1 & 0/10 \\
\bottomrule
\end{tabular}
\end{center}

Four independently swept architectures, four correlations
indistinguishable from zero. We report this as the paper's main
result.

We chose tables over error bars on the figure for the numeric
comparisons because Table~V must carry four quantities per row
(mean, $\sigma$, interval, F1) for ten rows; error bars would
convey the overlap qualitatively but not the interval widths, and
the tabular form costs less space under the six-page limit. Fig.~1
does carry confidence intervals, for the architecture means, where
a graphical comparison is the more useful presentation.

\subsection*{Point 3: Additional dataset}

\begin{rev}
(Consider) Test top 5--10 pairings on an additional sentiment
dataset (SST-2 / Amazon) to confirm the LSTM/GRU inversion holds
across domains.
\end{rev}

\noindent\pt{Not carried out. Named as the specific next step in
the Conclusion.}

We did not have the budget to add a second corpus within the
camera-ready window, having spent it on the 2{,}000-run
replication. We are not promising it here; we state plainly that
it is untested.

One clarification is owed to the reviewer, because the premise of
the suggestion no longer holds. The cross-architecture
\emph{inversion} of activation rankings that this point asks us to
confirm is not a real effect. It was an artefact of single-seed
measurement. On replicated data the six pairwise correlations
between architectures' activation marginals range from $-0.164$ to
$+0.576$ and none is significant ($p\geq0.082$): the rankings
neither transfer nor invert. There is consequently nothing to
confirm across domains on that specific point, and the claim has
been removed from the paper.

What a second corpus would genuinely test is whether the variance
split and the architectural advantage generalise, and the
Conclusion now names SST-2 and Amazon Reviews for exactly that
purpose. We agree it is the right next step.

\section*{Additional corrections made without prompting}

\begin{enumerate}[leftmargin=1.6em,itemsep=2pt]
  \item \pt{Fig.~1 caption was wrong.} It read ``Per-Epoch
  Validation Accuracy'' while the plot showed validation
  \emph{loss}. Corrected.

  \item \pt{A contradiction over seeds is resolved.} The submitted
  Conclusion attributed two results (88.48\%, 88.43\%) to ``a
  different seed in the convergence rerun'', while Table~VII stated
  the seed was fixed at 42. Auditing the code showed the opposite
  of what the text implied: the convergence rerun \emph{was} seeded
  at 42, and the main 400-run sweep set \emph{no seed at all}. Both
  statements have been removed. Over 2{,}000 seeded runs no
  configuration's mean reaches the 88.40\% benchmark and no
  individual run attained it (maximum 88.39\%), so the claim that
  the benchmark ``sits within normal seed variance'' is not
  supported and has been withdrawn.

  \item \pt{Table~IV row labels were corrupted} in the submitted
  source (``BiLSTM+BiLSTM+LSTM'', ``BiGRU+BiGRU+GRU''). Corrected.

  \item \pt{Claims withdrawn as unsupported by the replicated
  data.} Listed here for transparency: the identity of the single
  best activation pair; the cross-architecture inversion of
  rankings; the per-configuration precision--recall imbalances of
  the former Table~VIII (the widest gap between mean precision and
  mean recall over any cell is 6.89 points, while individual runs
  reach 24.17, so these are properties of runs, not of
  configurations); and the dead-neuron penalty for dense-layer
  ReLU of the former Table~X (no run of 2{,}000 fell below 60\%
  accuracy, and in BiLSTM+LSTM ReLU is in fact the strongest
  Act$_1$ choice, $p=0.006$). The corresponding tables have been
  removed and the space used for the variance and reproducibility
  results.
\end{enumerate}

\vspace{0.5em}
\noindent\rule{\textwidth}{0.4pt}

The paper is within the six-page limit, uses the IEEE conference
template, carries the required first-page header and copyright
footer, and has all fonts embedded. The single-seed observation was
the most useful comment we received; acting on it in full produced
a stronger result than the one we submitted.

\end{document}
